% !TEX root = ./main.tex
\documentclass[a4paper]{article}
% Course language: use [english] or [german].
\usepackage[german]{ajnotes}
\title{\textbf{Betriebwirtschaftslehre — Woche 3 Teil 2 Geschäftsplan}}

\begin{document}
\maketitle
\textit{siehe Folien} \url{https://lms.uzh.ch/auth/RepositoryEntry/17910989063/CourseNode/82289122621808}
    % Attach the given problem sheet: drop it in this folder as sheet.pdf and
    % un-comment the next line (pdfpages is already loaded by the preamble).
    % \includepdf[pages=-]{sheet.pdf}
    \begin{definition}
        [Geschäftsplan] definiert, welche Ziele in welchem Zeitraum mit welchen Mitteln erreicht werden sollen.
    \end{definition}
    \begin{itemize}
        \item Für Gründer 
        \subitem Systematisierung der Idee
        \subitem Zusammentraggen von Infos
        \subitem Vorbereiten der Entscheidungen
        \item Für Investor
        \subitem Abbau der Informationsasymmetrie zwischen Gründer und Investor 
        \subitem Abschätzung der Chancen Risiken der Investition
        \subitem Unternehmensbewertung (Valuation) und Rendite prognose
    \end{itemize}

    \section{Marketing}
    \begin{definition}
        [Marketing] Planung, Koordination und Kontrolle aller auf die aktuellen und potentiellen Märkte ausgerichteten Unternehmensaktivitäten.
    \end{definition}

    \subsection{Funktionen und Kern elemente}
    \begin{itemize}
        \item Analyse funktion der Inforamtionssammlung
        \begin{itemize}
            \item Gesamtmarkt abgrenzen \& Wachstums Abschätzung
        \item Wettbewerbs
        \end{itemize} 
        \item Entscheidungsfunktion 
        \begin{itemize}
        \item Wahl Zielmarkt
        \item Segmentierung und Wahl der Nische
        \item Differenzierung des Produkts (Positioning)
            
        \end{itemize}
        \item Umsetzungsfunktion (Unterstützung der Implementierung)
        \begin{itemize}
              \item 4Ps Product, Price, Place, Promotion
        \end{itemize}
        \item Ueberzeugen der Investoren
    \end{itemize}
    \subsection{Beispiel Madecasse}
    \begin{itemize}
        \item Kunden
        \item hohe Machtgegenüber madecasse wegen niedrige switching costs
        \item viele optionene 
        \item Lieferanten
        \item Hohe Verhandlungsmacht weil andere bereits etabliertfirmen höhere preise anbieten können
        \item kein Exklusiv vertrag
        \item Ersatzprodukte
        \item viele weil süssigkeiten und andere schokolade produkten sehr
        \item Potentielle Mitbewerber
        \item Hohe eintrittsbarrieren, bereits etablierte Konkurrenten
    \end{itemize}

    \subsection{Auswahl Zielmarkt}
    \subsubsection{Ziel}
    Übereinstimmung der eigenen Produktion mit Konsumenten markt, sodass Gesamtmarkt in kleineren \textit{homogenen} Bedürfnissesegmenten unterteilt werden. (intern homogen, extern heterogen)

    \subsubsection{Kriterien}
    \begin{itemize}
        \item GEographisch
        \item Demographisch
        \item Verhalten
        \item Lifestyle
        \item Einkaufsgewohnheit
    \end{itemize}
    \subsubsection{Beispiel Papierwindeln}
    \begin{itemize}
        \item Basis: 9.1 Mio schweizerbevölkerung
        \item Annahme: Avg Kind 2 jahre windelntragen (recherche)
        \item Basis: Lebenserwartung = ca 84 jahre
        \item Berechnung: 
    \end{itemize}

    \subsection{Gelstaltung der Organisation}
    \begin{itemize}
        \item Formell
        \begin{enumerate}
            \item interne Arbeitsteilung
            \item Koordinationsmechanismen
            \item Lohn und Anreizsystemen 
            \item Aus arbeitsteilung anforderungen
            \item Fähigkeitsanalyse
            \item Personalplanung
        \end{enumerate}
        \item Informell
        \begin{enumerate}
            \item Werte, Normen, Richtlinien
            \item Führungsstil etc
        \end{enumerate}
    \end{itemize}

    \subsection{Rechtsformen}
    \textbf{WICHTIG}
    \begin{itemize}
        \item Rechtsgemeinschaften (Unternehmer trägen direkte Verantwortung)
        \item Körperschaften (juristische Person trägt selbständige Verantwortung unabhängig der Mitglieder)
        \item 
    \end{itemize}
\end{document}
 